\ncdlesson{2}{Hvem er du, og hvor kommer du fra?}{你是谁，来自哪里？}
{DU3 Modul 1（A1）｜对标 NCE 1 Lesson 5--8}
{\textit{hv-} 疑问句倒装（\textit{hvad, hvor, hvem, hvordan, hvilke}）· 规则动词现在时（词干+\textit{r}）· 国籍、国家与语言名词 · \textit{hv-} 词首 \textit{h} 静音与常见弱读}
{本课是口语“信息交换”的核心地基。请先把 \textit{hvad, hvor, hvem, hvordan, hvilke} 五大疑问词读准（牢记 \textbf{h 完全不发音}），再通过第五部分把“国家大写、国籍语言小写且不加冠词”练成习惯。}

% ---------------------------------------------------------------------
\begin{foerduleaeser}
在阅读下面的迎新对话前，先看这两个问题，带着目标读课文：
\begin{enumerate}
  \item \textit{Hvor kommer Emma fra, og hvilke sprog taler hun?}（艾玛来自哪里？她说哪些语言？）
  \item \textit{Hvem er Marco, og hvilke dage har de undervisning?}（马尔科是谁？他们每周哪几天上课？）
\end{enumerate}
\end{foerduleaeser}

% ---------------------------------------------------------------------
\ncdsubsec{二、Dialog \& Udtale（课文对话、发音标注与中文译文）}

\begin{tipbox}{场景说明}
Emma 抵达奥胡斯语言中心（Aarhus Sprogcenter）办理新生迎新登记，前台接待员 Mette 为她录入个人信息；此时另一位新生 Marco 也走过来打招呼。
\end{tipbox}

\dialogturn{1}{Mette}{Godmorgen og velkommen til Aarhus Sprogcenter! Hvad hedder du?}
{\textit{godmorgen} \ipa{ɡoˈmɒːn}（\textbf{d 和 r 不单独弹出}，连读接近“哥莫恩”）；\textit{hvad} \ipa{ˈvað}，在快问句 \textit{hvad hedder du} 中常弱读为 \ipa{va ˈheðɐ du}（\textbf{h 和 d 都脱落}）。}
{早上好，欢迎来到奥胡斯语言中心！你叫什么名字？}

\dialogturn{2}{Emma}{Godmorgen! Jeg hedder Emma Wang. Undskyld, hvem er min lærer i dag?}
{\textit{hvem} \ipa{ˈvɛmˀ}（词首 \textbf{h 不发音}，带喉塞音）；\textit{lærer} \ipa{ˈlɛːɐ}（词尾 \textit{-er} 弱化成央低音 \ipa{ɐ}，听起来像长元音直接滑向“啊”）。}
{早上好！我叫王艾玛。请问，今天谁是我的老师？}

\dialogturn{3}{Mette}{Din lærer hedder Kasper. Hvordan staver du dit efternavn, Emma?}
{\textit{hvordan} \ipa{vɒˈdan}（\textbf{h 不发音}，重音在第二音节）；\textit{staver} \ipa{ˈsdæːvɐ}；\textit{efternavn} \ipa{ˈɛfdɐˌnɑʊ̯ˀn}（注意 \textit{-av-} 读作双元音 \ipa{ɑʊ̯}）。}
{你的老师叫卡斯珀。艾玛，你的姓氏怎么拼写？}

\dialogturn{4}{Emma}{Jeg staver det W-A-N-G. Jeg er ny kursist her på modul 1.}
{\textit{kursist} \ipa{kʰuɐ̯ˈsisd}（重音在第二音节 \textit{-sist}）；\textit{her} \ipa{ˈheɐ̯ˀ}（短促清晰带喉塞音）。}
{我拼作 W-A-N-G。我是这里 Modul 1 的新生。}

\dialogturn{5}{Mette}{Dejligt! Hvor kommer du fra, og hvilke sprog taler du?}
{\textit{dejligt} \ipa{ˈdɑɪ̯li}（\textbf{g 和 t 均不发音}！\textit{ej} 读 \ipa{ɑɪ̯}）；\textit{hvilke} \ipa{ˈvelɡə}（\textbf{h 不发音}，\textit{k} 读不送气清塞音 \ipa{ɡ}）；\textit{sprog} \ipa{ˈsbʁɔʊ̯ˀ}（\textit{-og} 读 \ipa{ɔʊ̯ˀ}）。}
{太好了！你从哪里来，会说哪些语言？}

\dialogturn{6}{Emma}{Jeg kommer fra Kina, så jeg er kineser. Jeg taler kinesisk, engelsk og lidt dansk.}
{\textit{kineser} \ipa{kʰiˈneːˀsɐ}、\textit{kinesisk} \ipa{kʰiˈneːˀsisg}（重音都在第二音节）；\textit{engelsk} \ipa{ˈɛŋˀəlsg}；\textit{lidt dansk} \ipa{led ˈdansg}（\textit{lidt} 末尾的 \textbf{d 不发音}）。}
{我来自中国，所以我是一名中国人。我说中文、英语和一点儿丹麦语。}

\dialogturn{7}{Mette}{Hvor bor du nu, og hvad laver du i Aarhus?}
{\textit{hvor bor du} \ipa{vɒː ˈboɐ̯ˀ du}（\textit{hvor} 弱读，重音落在 \textit{bor} 上）；\textit{hvad laver du} \ipa{va ˈlæːvɐ du}。}
{你现在住在哪里，在奥胡斯做什么工作？}

\dialogturn{8}{Emma}{Jeg bor i Aarhus, og jeg arbejder som ingeniør i et firma.}
{\textit{arbejder} \ipa{ˈɑːˌbɑɪ̯ˀdɐ}（第一音节 \textit{ar-} 读长音 \ipa{ɑː}，\textit{-ej-} 读 \ipa{ɑɪ̯}）；\textit{ingeniør} \ipa{eŋɕəˈnøɐ̯ˀ}（法语借词，\textit{g} 读如“西”的 \ipa{ɕ}，重音在最后）。}
{我住在奥胡斯，在一家公司当工程师。}

\dialogturn{9}{Marco}{Hej! Jeg hedder Marco. Jeg kommer fra Italien, og jeg lærer også dansk her. Hvilke dage har vi undervisning?}
{\textit{Italien} \ipa{iˈtʰæːljən}；\textit{dage} \ipa{ˈdæːə}（\textbf{中间 g 完全不发音}！）；\textit{undervisning} \ipa{ˈʌnɐˌviːˀsneŋ}（\textit{under-} 中 \textbf{d 不发音}）。}
{嗨！我叫马尔科。我来自意大利，我也在这里学丹麦语。我们哪几天有课？}

\dialogturn{10}{Mette}{I har undervisning tirsdag og torsdag. Velkommen på holdet, begge to!}
{\textit{tirsdag} \ipa{ˈtʰiɐ̯ˀsda}、\textit{torsdag} \ipa{ˈtʰɒːˀsda}（星期词尾 \textit{-dag} 的 \textbf{g 不发音}，弱读为 \ipa{da}）；\textit{begge to} \ipa{ˈbɛɡə ˈtʰoːˀ}（\textit{holdet} 读 \ipa{ˈhʌlˀəð}，\textbf{d 静音}）。}
{你们周二和周四上课。欢迎你们两位加入这个班！}

% ---------------------------------------------------------------------
\ncdsubsec{三、Gloser（核心生词表）}

\begin{tipbox}{词汇提示}
本课重点掌握 5 个 \textit{hv-} 疑问词、6 个高频规则动词的现在时，以及国家、国籍与语言的词形规律。
\end{tipbox}

\begin{glosertable}
疑问词 & hvad / hvem & --- & 什么 / 谁 \\
疑问词 & hvor / hvordan & --- & 哪里 / 怎么样，如何 \\
疑问词 & hvilken & \textit{(hvilket, hvilke)} & 哪一个 / 哪些（随名词性数变化） \\
\midrule
名词 & en lærer & \textit{(læreren, lærere, lærerne)} & 老师，教师 \\
名词 & en kursist & \textit{(kursisten, kursister, kursisterne)} & 语言学校学员 \\
名词 & et efternavn & \textit{(efternavnet, efternavne, efternavnene)} & 姓氏 \\
名词 & et sprog & \textit{(sproget, sprog, sprogene)} & 语言（单复数同形） \\
名词 & et hold & \textit{(holdet, hold, holdene)} & 班级；团队（单复数同形） \\
名词 & en undervisning & \textit{(undervisningen, 不可数)} & 教学，课程 \\
名词 & et firma & \textit{(firmaet, firmaer, firmaerne)} & 公司，企业 \\
\midrule
动词 & at stave & \textit{(staver, stavede, har stavet)} & 拼写 \\
动词 & at tale & \textit{(taler, talte, har talt)} & 说（语言）；交谈 \\
动词 & at bo & \textit{(bor, boede, har boet)} & 居住 \\
动词 & at lave & \textit{(laver, lavede, har lavet)} & 做，制作（\textit{hvad laver du?} 你干嘛/从事什么） \\
动词 & at arbejde & \textit{(arbejder, arbejdede, har arbejdet)} & 工作 \\
动词 & at lære & \textit{(lærer, lærte, har lært)} & 学习；学会 \\
\midrule
形容词 & ny & \textit{(nyt, nye)} & 新的 \\
副词 & lidt & --- & 一点儿（\textit{lidt dansk} 一点丹麦语） \\
\end{glosertable}

% ---------------------------------------------------------------------
\ncdsubsec{四、Grammatik \& Faste vendinger（核心语法与地道表达精讲）}

\begin{grammatikbox}{语法点 1：\textit{hv-} 特殊疑问句（Hv-spørgsmål）与倒装语序}
丹麦语所有特殊疑问词都以 \dabold{hv-} 开头，且\textbf{词首字母 h 永远不发音}。疑问句同样遵循\textbf{V2 原则}——\dabold{变位动词永远占第 2 位}：
\begin{quote}
\textbf{公式}：\dabold{Hv-疑问词} + \dabold{动词} + \da{主语} + 其他成分？（若 \textit{hvem/hvad} 本身就是主语，则直接接动词）
\end{quote}

\vspace{0.2em}
\renewcommand{\arraystretch}{1.25}
{\small
\begin{tabularx}{\linewidth}{@{} >{\bfseries\color{NCDNavy}}l >{\raggedright\arraybackslash}p{3.2cm} >{\raggedright\arraybackslash}X @{}}
\toprule
\textbf{疑问词} & \textbf{基本含义} & \textbf{课文与实战例句（注意动词在第 2 位）} \\
\midrule
hvad \ipa{ˈvað} & 什么 & \textit{\dabold{Hvad} hedder du?} / \textit{\dabold{Hvad} laver du i Aarhus?} \\
hvem \ipa{ˈvɛmˀ} & 谁 & \textit{\dabold{Hvem} er min lærer?}（\textit{hvem} 作主语） \\
hvor \ipa{ˈvɒːˀ} & 哪里（+形容词：多么） & \textit{\dabold{Hvor} kommer du fra?} / \textit{\dabold{Hvor} bor du nu?} \\
hvordan \ipa{vɒˈdan} & 怎样，如何 & \textit{\dabold{Hvordan} staver du dit efternavn?} \\
hvilken / -t / -e & 哪一个 / 哪些（+名词） & \textit{\dabold{Hvilke sprog} taler du?} / \textit{\dabold{Hvilke dage} har vi undervisning?} \\
\bottomrule
\end{tabularx}
}
\vspace{0.2em}

\textbf{注意介词后置}：问“来自哪里”时，\textit{fra} 不放句首，而是放到句尾：\textit{\textbf{Hvor} kommer du \dabold{fra}?}（不要说 *\textit{Fra hvor kommer du?}）。
\end{grammatikbox}

\begin{grammatikbox}{语法点 2：规则动词现在时（Nutid）构成法：词干 + \textit{r}}
丹麦语绝大多数动词的现在时极其简单——\textbf{在不定式（\textit{at} 后的原形）词尾直接加 \dabold{-r}}，且全部人称通用：
\begin{itemize}
  \item \textbf{以非重读 \textit{-e} 结尾的不定式}（最常见）：直接加 \textit{-r} 变成 \textit{-er}。\newline 例：\textit{at tale} $\rightarrow$ \da{taler}，\textit{at stave} $\rightarrow$ \da{staver}，\textit{at arbejde} $\rightarrow$ \da{arbejder}，\textit{at lære} $\rightarrow$ \da{lærer}，\textit{at komme} $\rightarrow$ \da{kommer}。
  \item \textbf{以重读元音结尾的不定式}：直接加 \textit{-r}。\newline 例：\textit{at bo} $\rightarrow$ \da{bor}，\textit{at gå} $\rightarrow$ \da{går}，\textit{at stå} $\rightarrow$ \da{står}，\textit{at se} $\rightarrow$ \da{ser}。
  \item \textbf{极少数需单独记忆的特殊现在时}：\textit{at være} $\rightarrow$ \dabold{er}（是），\textit{at have} $\rightarrow$ \dabold{har}（有），\textit{at vide} $\rightarrow$ \dabold{ved}（知道），\textit{at gøre} $\rightarrow$ \dabold{gør}（做）。
\end{itemize}
\textbf{语用提示}：丹麦语没有专门的“现在进行时”形式，\textit{Jeg arbejder i Aarhus} 既可表示“我在奥胡斯工作（长期）”，也可表示“我正在工作（此刻）”。
\end{grammatikbox}

\begin{grammatikbox}{语法点 3：国家（Land）、国籍（Nationalitet）与语言（Sprog）三大铁律}
丹麦语的大小写与冠词规则和英语正好相反，考场极易扣分：

\vspace{0.25em}
\renewcommand{\arraystretch}{1.22}
{\small
\begin{tabularx}{\linewidth}{@{} >{\raggedright\arraybackslash}p{3.0cm} >{\raggedright\arraybackslash}p{3.1cm} >{\raggedright\arraybackslash}p{2.7cm} >{\raggedright\arraybackslash}X @{}}
\toprule
\textbf{国家（\dabold{大写}）} & \textbf{国籍名词（\dabold{小写}）} & \textbf{语言 / 形容词（\dabold{小写}）} & \textbf{自我介绍例句} \\
\midrule
\da{D}anmark & en \da{d}ansker & \da{d}ansk & \textit{Han er dansker og taler dansk.} \\
\da{K}ina & en \da{k}ineser & \da{k}inesisk & \textit{Jeg er kineser og taler kinesisk.} \\
\da{I}talien & en \da{i}taliener & \da{i}taliensk & \textit{Marco er italiener og taler italiensk.} \\
\da{T}yskland & en \da{t}ysker & \da{t}ysk & \textit{Hun er tysker og taler tysk.} \\
\da{S}torbritannien / \da{U}SA & en \da{b}rite / \da{a}merikaner & \da{e}ngelsk & \textit{Vi taler engelsk sammen.} \\
\bottomrule
\end{tabularx}
}
\vspace{0.25em}

\begin{itemize}
  \item \textbf{铁律 \ding{172} 大小写}：只有\textbf{国家名}首字母大写（\textit{Kina, Danmark}）；\textbf{国籍和语言全部小写}（\textit{kineser, dansk, engelsk}）。
  \item \textbf{铁律 \ding{173} 零冠词}：用 \textit{være} 介绍\textbf{国籍或职业}时，\textbf{绝对不加不定冠词 \textit{en}}！\newline 对：\textit{Jeg er \da{kineser}.} / \textit{Jeg er \da{ingeniør}.} \quad 错：*\textit{Jeg er en kineser / en ingeniør.}
\end{itemize}
\end{grammatikbox}

\begin{grammatikbox}{语法点 4：地道交际表达（Faste vendinger）与 \textit{hv-} 静音小结}
\begin{itemize}
  \item \da{Hvad laver du?} —— 视语境既可问“你现在在干嘛？”，也可问“你做什么工作？”。
  \item \da{Hvordan staver du det?} —— 你怎么拼写它？（口语与填表必用句）。
  \item \da{Velkommen til [地点] / Velkommen på [班级/团队]} —— 欢迎来到…… / 欢迎加入……。
  \item \da{begge to} —— 两个人一起，俩人都（比只说 \textit{begge} 更地道口语化）。
  \item \textbf{发音口诀}：凡见 \textit{hv-} 开头（\textit{hvad, hvor, hvem, hvordan, hvilken, hvornår}），只读 \ipa{v} 不读 \ipa{h}；同时 \textit{h} 在 \textit{hj-}（如 \textit{hjem} \ipa{ˈjɛmˀ}）中也不发音！
\end{itemize}
\end{grammatikbox}

% ---------------------------------------------------------------------
\ncdsubsec{五、Mønsterøvelser（新概念式句型操练）}

\textbf{练习 A：用正确的 \textit{hv-} 疑问词填空（\textit{hvad, hvor, hvem, hvordan, hvilke}）}
\begin{enumerate}
  \item \underline{\hspace{2.2cm}} hedder din lærer på sprogskolen?
  \item \underline{\hspace{2.2cm}} bor du i Danmark, og \underline{\hspace{2.2cm}} kommer du fra?
  \item \underline{\hspace{2.2cm}} staver du dit fornavn?
  \item \underline{\hspace{2.2cm}} sprog taler I på holdet?
  \item \underline{\hspace{2.2cm}} er ny kursist på modul 1? --- Det er Emma og Marco.
\end{enumerate}

\textbf{练习 B：将括号内的动词变为现在时，注意国籍与职业的零冠词规则}
\begin{enumerate}
  \item Emma \underline{\hspace{2.0cm}} (\textit{at komme}) fra Kina. Hun \underline{\hspace{2.0cm}} (\textit{at være}) kineser.
  \item Nu \underline{\hspace{2.0cm}} (\textit{at bo}) hun i Aarhus og \underline{\hspace{2.0cm}} (\textit{at arbejde}) som ingeniør.
  \item Marco \underline{\hspace{2.0cm}} (\textit{at lære}) dansk, og han \underline{\hspace{2.0cm}} (\textit{at tale}) allerede lidt dansk.
  \item Hvilke dage \underline{\hspace{2.0cm}} (\textit{at have}) I undervisning?
  \item Jeg \underline{\hspace{2.0cm}} (\textit{at stave}) mit efternavn på dansk.
\end{enumerate}

\textbf{练习 C：中译丹实战（注意 \textit{hv-} 语序、大小写及零冠词）}
\begin{enumerate}
  \item 你叫什么名字，你来自哪个国家？（提示：哪个国家 \textit{hvilket land}）
  \item 我是一名中国人，但我住在丹麦。
  \item 你说哪些语言？--- 我说中文、英语和一点儿丹麦语。
  \item 他在公司当老师，每周二和周四上课。
  \item 请问，你现在在奥胡斯做什么工作？
\end{enumerate}

% ---------------------------------------------------------------------
\ncdsubsec{六、Resumé \& Mundtlig træning（口语复述与参考答案）}

\textbf{口语复述任务}：遮住课文，先口头回答 \textbf{Før du læser} 的两个问题，然后用 \textbf{4--5 句话} 分别介绍你自己（姓名、国籍、语言、居住城市与职业/学习情况）以及课文的 Emma 和 Marco。

\begin{answerbox}
\textbf{练习 A 答案}
\begin{enumerate}
  \item \da{Hvad}（问名字固定用 \textit{Hvad hedder...?}，不要受英语“Who/How”影响）
  \item \da{Hvor} ... \da{hvor}（问居住地与来源地均用 \textit{hvor}，后者配合句尾 \textit{fra}）
  \item \da{Hvordan}（问拼写方式“如何拼”）
  \item \da{Hvilke}（\textit{sprog} 在此指复数“哪些语言”，配 \textit{hvilke}）
  \item \da{Hvem}（从回答 \textit{Det er Emma og Marco} 可知问人，用 \textit{hvem}）
\end{enumerate}

\textbf{练习 B 答案}
\begin{enumerate}
  \item \da{kommer}；\da{er}（\textit{at være} 现在时不规则变为 \textit{er}）
  \item \da{bor}（元音结尾直接加 \textit{-r}）；\da{arbejder}
  \item \da{lærer}；\da{taler}
  \item \da{har}（\textit{at have} 现在时变为 \textit{har}）
  \item \da{staver}
\end{enumerate}

\textbf{练习 C 答案}
\begin{enumerate}
  \item \textit{Hvad hedder du, og hvilket land kommer du fra?}
  \item \textit{Jeg er kineser, men jeg bor i Danmark.}（\textit{kineser} 小写且无 \textit{en}，\textit{Danmark} 大写）
  \item \textit{Hvilke sprog taler du? --- Jeg taler kinesisk, engelsk og lidt dansk.}（语言全部小写）
  \item \textit{Han arbejder som lærer i et firma, og han har undervisning tirsdag og torsdag.}
  \item \textit{Undskyld, hvad laver du i Aarhus nu?}
\end{enumerate}

\textbf{Før du læser 参考回答 \& 口语复述范文（Resumé）}
\begin{itemize}
  \item \textbf{问题 1}：\textit{Emma kommer fra Kina. Hun taler kinesisk, engelsk og lidt dansk.}
  \item \textbf{问题 2}：\textit{Marco er en ny kursist fra Italien. De har undervisning tirsdag og torsdag.}
  \item \textbf{口语复述范文}：\textit{Emma er ny kursist på Aarhus Sprogcenter. Hun kommer fra Kina og bor nu i Aarhus, hvor hun arbejder som ingeniør. Hun taler kinesisk, engelsk og lidt dansk. Hendes lærer hedder Kasper, og en anden kursist på holdet hedder Marco og kommer fra Italien. De har undervisning tirsdag og torsdag.}
\end{itemize}
\end{answerbox}
